% paper/sections/abstract.tex — revised in the pre-arXiv review pass.
Functional data analysis treats curves, surfaces, and trajectories as the
primary unit of observation and provides a rich family of methods---smoothing
and basis representation, registration, depth and outlier detection, functional
principal component analysis (FPCA), clustering, classification, regression,
forecasting, and process monitoring.
\texttt{fdars} is a Python library that brings this breadth together in a
single package: \nsubmodules{} analysis submodules exposing
\npubliccallables{} public callables (plus the advisor, scikit-learn, plotting,
and MCP layers), implemented in a Rust core (\texttt{fdars-core}) and exposed through
PyO3 bindings.
A scikit-learn-compatible layer provides \nsklearnestimators{} estimators that
pass \texttt{check\_estimator}, so functional methods compose with scikit-learn
pipelines and cross-validation.
Two further components were not found in any of the packages we compared: a
grounded parameter advisor (\texttt{fdars.advisor}) that reasons over
deterministic, library-computed diagnostics and rejects any advice citing a
numeric value not present in them, and a machine-readable
scientific-provenance map that links callables to the publications
defining them (currently \ncoverage{} callables with curated, author-verified
entries).
We describe the architecture and data model, compare coverage with Python, R,
and MATLAB software, check numerical agreement against independent reference
implementations, and illustrate the library in four case studies on real data.
Every code snippet, figure, case-study result, and cross-implementation
comparison in this paper is produced by committed, executed scripts;
continuous integration checks the snippets, case-study numbers, and the
\texttt{fdars} side of the agreement table for drift, and checks that figure
generation is deterministic.
